\section{Transition d'un framework RAC vers LTS : PDDL}\label{parser_section}

Les LTS ne constitue pas directement un langage pour le domaine du Raisonnement sur l'Action et le changement (RAC) mais, comme dit dans \cite{sarmiento_formalising_2025} leur utilisation permet d'unifier des raisonnements sur différents frameworks ainsi que de lever certaines ambiguïtés. Tout formalisme pouvant être représenté sous forme de LTS bénéficie ainsi des réflexions et observation de ce travail. Pour illustrer la possibilité de transition, nous allons développer sur la traduction d'un problème PDDL vers LTS. La représentation LTS permet de travailler sur la comparaison de contexte causal contrefactuel, mais on laisse le travail de recherche et de résolution de problème de planification aux outils préexistants. En effet, la complexité exponentielle d'un LTS en fonction du nombre d'états et d'action n'est pas adapté à la recherche de politique.


\subsection{Rappel PDDL}

Un problème de planification est décrit en PDDL par deux types de fichiers. Le premier définit le domaine de planification, il contient la déclaration des types, des actions (effets et préconditions) ainsi que les axiomes et les règles du domaine. Le second fichier définit le problème, il indique les objets, les constantes spécifiques ainsi que la situation de départ et l'objectif. Un exemple se trouve en annexe \ref{ex_pddl}


\subsection{Construction d'un LTS à partir d'un problème PDDL}
Pour construire un LTS à partir d'un couple de fichier domaine/problème, on commence par instancier toutes les actions possibles à partir de leurs définitions et des objets du problème, ce processus est explicité dans ci-dessous.

\hfill
\begin{lstlisting}[language=Python,basicstyle=\ttfamily\footnotesize,
                   commentstyle=\color{gray}\itshape,
                   keywordstyle=\bfseries]
def instantiate(actions, objects):
    instanciated = []
    for action in actions:
        params = action.parameters
        for args in product(objects, repeat=len(params)):
            val = dict(zip(params, args))
            name, precond,effect =  action.apply(val)  
            instanciated.append((name, precond, effect))
    return instanciated
\end{lstlisting}
\hfill

Une fois l'ensemble de toutes les actions instanciées, avant de pouvoir construire le LTS, on définit la fonction de satisfaction de préconditions. Pour un ensemble de préconditions et un état représenté par un ensemble de fluents, elle vérifie si les préconditions sont respectées dans l'état. Cette fonction gère les préconditions négatives, disjonctives et conjonctives.

\hfill
\begin{lstlisting}[language=Python,basicstyle=\ttfamily\footnotesize,
                   commentstyle=\color{gray}\itshape,
                   keywordstyle=\bfseries]
def statisfied(preconditions,state):
    if preconditions conjonctives:
        return all(statisfied(p,state) for p in preconditions)
    if preconditions negatives:
        return not statisfied(preconditions,state)
    if preconditions disjonctives:
        return any(statisfied(p,state) for p in preconditions)
    return precondition in state

\end{lstlisting}
\hfill

La seconde fonction qui est nécessaire avant celle de construction du graphe est la fonction d'application des effets positifs et négatifs d'une action dans un état.
\hfill

\begin{lstlisting}[language=Python, basicstyle=\small\ttfamily,
                   commentstyle=\color{gray}\itshape,
                   keywordstyle=\bfseries]
def apply_effects(effects,state):
    new_state = state
    for f in effects:
        if isinstance(f, Not):
            new_state.discard(f) # effet negatif
        else:
            new_state.add(f) #effet positif
    return new_state

\end{lstlisting}
\hfill

Avec ces trois fonctions, on peut maintenant définir la fonction de construction du LTS associé au problème PDDL. Un état est défini par l'ensemble des prédicats qui sont vrais dans ce dernier. On utilise une pile pour progresser dans la construction du graphe, à chaque nouvel état, on observe parmi toutes les actions possibles lesquelles sont applicables en fonction de leurs préconditions. Si une action est applicable, on ajoute une nouvelle arête qui a pour état d'arrivée soit un état préexistant soit un nouvel état que l'on ajoute à la pile. Les états finaux sont les états dans lesquelles les conditions de l'objectif sont respectés. Pour les distinguer des autres états une arête réflexive est ajoutée. 

\hfill
\begin{lstlisting}[language=Python,basicstyle=\ttfamily\footnotesize,
                   commentstyle=\color{gray}\itshape,
                   keywordstyle=\bfseries]
def build_graph(actions, init_state, goals):
    state = init_state
    G = graph
    G.add_node(state)
    queue = [state]
    visited = {}
    while queue:
        current_node = queue.pop(0)
        if current_node in visited:
            continue
        visited.add(current_node)

        for preconditions, effects in actions:
            if satisfied(preconditions, current_node):
                new_node = apply_effect(effects, current_node)

            exist = new_node in G.nodes
            if not exist:
                G.add_node(new_node)
            else:
                continue
            
            G.add_edge(current_node, new_node)

            if new_node not in visited:
                queue.append(visited)
    
    for node in G.nodes:
        if satisfied(goals, node):
            G.add_edge(node,node, action = '$\\varnothing')

\end{lstlisting}
\hfill



\subsection{Remarques et limites}

% Les LTS obtenus ne représentent pas l'ensemble des états possibles du monde, seuls les états accessibles depuis l'état initial, respectant les fonctions de précondition y figurent. Pour enrichir ces graphes, il serait potentiellement intéressant de construire pour un même domaine et objectif,  différents graphes selon différents états de départ et les fusionner afin de pouvoir travailler notamment sur des contextes contrefactuels avec des états de départ qui diffèrent (peut être potentiellement très coûteux). 
Il est possible de continuer de tirer parti de PDDL et de ses solveurs pour extraire de manière automatique des politiques qui nous permettent par la suite de travailler sur des contextes causaux. Pour construire la politique, on avance dans le graphe depuis l'état de départ en attribuant à chaque état dans lequel on se trouve l'action associée dans le plan. Si besoin ces politiques partielles peuvent ensuite être complétées aléatoirement ou selon une heuristique.

Il est à noter que les LTS correspondent à une sémantique et non pas à de vrais outils de travail, leur utilité pour des applications concrètes reste limitée et ils seront surtout utilisés de manière formelle pour raisonner. Pour des applications à des exemples simples, ils permettent une meilleure visualisation, on considère donc suffisant un parseur simple comme présenté ci-dessus. L'idée de cette section étant principalement d'expliciter la faisabilité d'une transition d'un framework RAC vers des LTS, de possibles améliorations sont laissées en suspend, car ne représentant pas le cœur du travail. On peut par exemple noter que ce parseur ne gère que des problèmes PDDL contenant uniquement le requirement STRIPS et negative-preconditions. 


